\documentclass{book}

\usepackage[T1]{fontenc}
\usepackage[portuguese]{babel}


\usepackage{manual_template}
% tem que decidir uma fonte legal
% https://tug.org/FontCatalogue/typewriterfonts.html

% para gerar texto lixo
\usepackage{blindtext}
\usepackage{lipsum}

\title{Manual de Ingressante}
\author{Centro Acadêmico da Computação}
\newcommand{\ano}{2024}
\date{\ano}
\newcommand{\shortauthor}{CACo}

\newcommand{\autoresCapa}{%
autor1 (EC021) \\
autor 2 (CC020)
}%

\newcommand{\autoresDiagramacao}{%
autora.e (curso) \\
autora.e (curso) \\
autora.e (curso)
}%

\newcommand{\autoresRevisao}{%
autora.e (curso) \\
autora.e (curso) \\
autora.e (curso) \\
autora.e (curso) \\
autora.e (curso) \\
autora.e (curso)
}%

\newcommand{\gestao}{nome da gestao (\ano)}

\begin{document}

%capa
.
\vfill
{\Huge Aqui é a capa}
\vfill
\thispagestyle{empty}\cleardoublepage

%contracapa
{\color{white}.}
\vfill
\begin{center}
    {\color{darkgray}\bfFont\fontsize{35}{35}\selectfont Manual de Ingressante}

    %imagem
    \vspace{3em}
    {imagem caquinho}
    \vspace{3em}
    %imagem

    \vfill
    {\color{darkgray}\generalFont\normalsize\selectfont Centro Acadêmico da Computação \\ Universidade Estadual de Campinas \\
    \vspace{1em}
    Fevereiro \ano
    \vspace{1em}
    \begin{tags} \instagram{cacounicamp} \sep \email{caco.gestao@gmail.com}\end{tags}
    }
    \vspace{3em}

\end{center}
\vfill
\thispagestyle{empty}\cleardoublepage

%licenca e autores
\begin{center}
    \begin{figure}[h]
        \centering
        \includegraphics[width=0.3\textwidth]{_adendos/CC-BY-SA.jpg}
    \end{figure}    

    \vfill
    {\color{darkgray}\generalFont\normalsize\selectfont
    
    Este trabalho está licenciado sob a Licença Attribution ShareAlike 3.0 Brasil da Creative Commons. Para ver uma cópia desta licença, visite: \href{https://creativecommons.org/licenses/by-sa/3.0/br}{creativecommons.org/licenses/by-sa/3.0/br}
    
    \vfill
    Manuais anteriores e seus textos corridos estão disponíveis em: \github{cacounicamp/Manual-do-Bixo}

    \vfill
    Capa e verso: \\
    \autoresCapa

    \vfill
    Diagramação: \\
    \autoresDiagramacao

    \vfill
    Revisão: \\
    \autoresRevisao

    \vfill
    Gestão \gestao

    \vfill
    Diagramado em LaTeX e ilustrado em Figma. \\
    \vfill
    Impresso na Gráfica Central da Unicamp em couchê 250g/m2 na capa e offset 75g/m2 no miolo.
    }
\end{center}
%\vfill
\thispagestyle{empty}\cleardoublepage


\clearpage                       % Otherwise \pagestyle affects the previous page.
{
  \pagestyle{empty}
  \fancypagestyle{plain}         % Re-definition removes numbers from first page.
  {
    \fancyhf{}%                       % Clear all header and footer fields.
    \renewcommand{\headrulewidth}{0pt}% Clear rules (remove these two lines if not desired).
    \renewcommand{\footrulewidth}{0pt}%
  }
  \tableofcontents
  \thispagestyle{empty}
}


%\maketitle
%\tableofcontents

%\setcounter{page}{1}

\chapter{Olá, Mundo!}

\setcounter{page}{1}

\section{Prefácio}

Bem-vinde!

Este é o Manual de Ingressante do CACo, o Centro Acadêmico da Computação. Esse manual foi elaborado por diverses veteranes e inclui várias dicas para auxiliar a sua sobrevivência nestes primeiros anos na Unicamp. Você deve receber a versão impressa dele nos primeiros dias ou semanas de aula! Se ainda não a pegou, entre em contato com alguém do CACo.

Nós trabalhamos muito para que as informações aqui presentes sejam úteis e atuais; e queremos ouvir de você seu feedback sobre este material! Se você encontrar quaisquer informações incorretas ou desatualizada, nos avise por e-mail ou por mensagem no Instagram.

\begin{tags}
    \email{caco.gestao@gmail.com} \sep \instagram{cacounicamp}    
\end{tags}

\subsection{Sobre o gênero dos substantivos neste manual}

Você, ex-vestibulande, humane do século XXI, certamente está antenade na lutas sociais que acontecem à sua volta e uma delas é a da igualdade de gêneros. Ela se baseia no respeito e inclusão de todos os gêneros, de forma que todes tenham os mesmos direitos em nossa sociedade - inclusive o direito a ser corretamente representade.

Infelizmente, nosso idioma não foi feito levando em conta a igualdade de gênero. Isso acabou gerando discrepâncias como a regra irracional e machista de que devemos nos referir a um grupo de 100 mulheres e apenas um homem utilizando o pronome "eles". Por esse motivo, utilizamos em nosso manual uma linguagem diferente, centrada nos artigos neutros.

Como estamos acostumados ao português padrão e este manual passa por constantes atualizações, talvez você encontre lugares onde o texto ainda não utiliza uma linguagem neutra. Se isso acontecer, nos envie uma solicitação para alterarmos o texto e nos ajude a deixar o manual cada vez melhor para todo mundo.

Vale enfatizar que fazemos isso porque valorizamos as minorias presentes em nossa comunidade e deixamos a seguinte mensagem para todes es seus representantes: \textbf{Nós, o seu Centro Acadêmico, estamos aqui para você!}

\section{Mensagem da FEEC}

<mensagem da FEEC>


\section{Mensagem do IC}

<mensagem do ic>

\section{Os melhores anos da sua vida!}

Parabéns!

Entrar na Universidade e cursar o ensino superior é uma grande conquista e pode lhe trazer algumas das experiências mais interessantes da sua vida. Uma delas é a oportunidade de se especializar em uma área do conhecimento e aprender com cientistas e profissionais que sabem muito bem - ou até mesmo descobriram - os assuntos que estão ensinando.

Nunca é tarde demais para estudar, mas muites de nós, quando temos essa oportunidade, a agarramos ainda jovens. E assim é comum que muitas coisas se misturem aos desafios da vida acadêmica: a entrada no mundo adulto, relações marcantes entre colegas e amigues, dezenas de convites para festas por mês, responsabilidades crescentes e as famosas contas pra pagar.

Enfim, o peso das novas responsabilidades e dos estudos, assim como as deficiências e exigências da própria universidade podem impactar nossa saúde mental - um tema que está sendo cada vez mais discutido e que incluímos neste manual para você poder começar o curso mais tranquile.

Mas muito mais do que estudar, aqui você vai encontrar toda uma comunidade realizando uma serie de atividades extras incríveis. Atividades para todos os gostos: de grupos de estudo a redes de apoio, times de esportes e e-sports, atléticas, empresas juniores e nós o seu centro acadêmico. Essas experiências elevam o significado dos anos na universidade a muito mais do que aulas, provas, livros e exercícios. E, através delas, você pode encontrar amizades pra toda a vida!

Entrar numa universidade internacionalmente reconhecida, como a Unicamp, tem ainda mais desdobramentos. Suas atividades cotidianas poderão ter um impacto tremendamente grande lá fora!

Primeiro, porque nossa universidade é pública. Isso significa que você é ume des pouques que a população brasileira conseguiu colocar pra dentro de um ensino superior de qualidade que não apenas ensina uma nova profissão, mas que também produz ciência. De fato, as universidades públicas produzem 95\% da ciência do nosso país \footnote{buscar a referência}. Não por acaso, a poucos quilômetros do campus da Unicamp fica o Sírius, um acelerador de partículas de padrão mundial! Por esses motivos, pelo nosso trabalho científico e de nossos colegas, podemos nos orgulhar do retorno que damos à sociedade, fazendo o peso do nome Unicamp valer muito mais do que um diploma na parede.

Segundo, porque as universidades públicas estão em disputa e sofrendo ataques, com constantes tentativas de descredibilização e precarização. Houve cortes de verbas bilionários que afetaram as universidades federais por motivos de "balbúrdia", de acordo com o ex-ministro da Educação, Abraham Weintraub. Recentemente, após a pandemia e o
aprofundamento da crise econômica, os ataques têm sido cada vez mais intensos e o ensino superior tem sofrido de uma crise orçamentária em níveis nacionais.

A Unicamp não é uma bolha: esses ataques tem consequências aqui. Sofremos cortes de verbas; há muita terceirização, diminuindo a qualidade do serviço e as condições de trabalho e sustento dos funcionários; faltam vagas na moradia estudantil e bolsas para estudantes mais pobres. A iniciativa privada age predatoriamente em prol desse desmonte e se infiltram cada vez mais fundo na universidade, controlando quais pesquisas serão realizadas ou não em prol dos seus próprios interesses - diminuindo seus gastos através do uso da infraestrutura paga com dinheiro público. Isso impede a democratização do conhecimento e prejudica inclusive quem pesquisa, que não só se torna refém como passa a ter uma empresa decidindo o que pode ou não ser publicado. Essas são questões fundamentais que você vai conhecer mais adiante neste manual e no
seu dia a dia.

A realidade é que a universidade pública precisa que você a defenda como parte de sua jornada por aqui. A ciência e a educação são grandes patrimônios do humanidade e somos nós os responsáveis por zelar por elas. Não desperdice a sua chance de cumprir esse papel!

Por tudo isso, bem-vinde! Os melhores anos da sua vida serão o que você fizer que eles sejam. Você tem um longo caminho a trilhar agora e vai se deparar com muitos desafios. Então, não se esqueça: você sempre pode contar com a gente.


\chapter{Sobrevivendo na Unicamp}

\section{A palavra é permanência}

A ideia é simples: agora que você entrou na universidade, não vá embora tão cedo e
sem o diploma. Você precisa permanecer. Isso envolve uma porção de fatores e
necessidades diferentes para cada estudante.

Ter o que comer e onde morar: permanência. Ter estabilidade emocional e saúde
mental: permanência. Ter como se locomover até a universidade e outros locais de
estudo: permanência. Ter acesso aos livros, conteúdos, softwares, computadores,
equipamentos, sistema de saúde: tudo isso é permanência.

Qualquer um desses pilares, se cair, pode te deixar numa situação extremamente
vulnerável e até mesmo forçar a sua evasão da graduação. Afinal de contas, todos esses
fatores impactam diretamente no seu desempenho acadêmico e na sua capacidade de
sobreviver enquanto cumpre o currículo. Mas como garantimos a nossa permanência?

\subsection{Dois modelos de universidade}

Um dos jeitos é o "se vira". Para estudantes com famílias mais bem estruturadas
economicamente, é razoavelmente simples garantir tudo isso: pais pagam o aluguel; o
bandejão tem um preço confortável mas você nem precisa dele; pode pagar os próprios
livros; você tem seu computador. E quem não tiver essas coisas garantidas que se vire e
procure um emprego. Parece coisa de universidade particular, não é? Um dos lugares
onde a política do "se vira" é a norma são países em que o Ensino Superior é todo
privatizado, como os EUA, onde os estudantes já chegaram a acumular US\$ 1.5 trilhão
em dívidas para pagar despesas com a faculdade\footnote{\textit{Student Loan Debt Statistics In 2019: A \$1.5 Trillion Crisis. Forbes, 2019.}}. A problemática é tão grave que há
estudantes morando dentro de carros para não perder a vaga em universidades
prestigiadas\footnote{\textit{Number of Homeless Students Soars. US News, 2019.}} – o que para alguns pode ser visto como mais uma história de
superação, é uma realidade dura e injusta para milhares de estudantes.

Colocando isso em perspectiva fica claro como a universidade pública pode
democratizar o conhecimento. Ainda assim, ela não garante automaticamente a
permanência de todes es seus alunes. Não muito tempo atrás, mesmo que você
conseguisse entrar em uma universidade pública como a Unicamp, apesar de não ter
que pagar mensalidade, ainda ia ser bem difícil se virar pra permanecer por aqui.

Assim, o outro jeito de garantir permanência é que a sociedade como um todo arque
com os custos do cumprimento do direito irrestrito e constitucional à Educação dos
estudantes através da realização de Políticas de Permanência. Mesmo que o
investimento seja alto, o retorno trazido por ume profissional de ponta à sociedade mais
do que paga esse investimento em um período de tempo razoavelmente curto.

É importante pontuar que as Políticas de Permanência que temos hoje tiveram de ser
conquistadas. Estudantes lutaram pesado, bravamente, para que hoje tenhamos direito a
Moradia Estudantil, ao Bandeijão, as bolsas de auxilio social, entre outras Políticas de
Permanência. Além delas, também foram conquistadas Políticas Afirmativas como a
bonificação para alunos de escolas públicas no vestibular, cotas étnico-raciais e o
vestibular indígena que fazem com que a universidade pública incorpore cada vez mais
a nossa sociedade. Como você vai ver no próximo texto, “Moradia: luta e disputa”, as
coisas ainda estão longe de serem perfeitas, essas lutas não pararam no tempo, ainda
temos muito o que conquistar e precisamos estar sempre atentes porque o que já
conquistamos pode ser tomado de nós.

<imagem unicamp>

\section{Moradia: luta e disputa}

\textit{Texto adaptado, atualizado e ampliado do site do Cêntro Acadêmico do IEL (CAL).}

Em 1986, foi organizado o movimento estudantil TABA: sem ter onde morar, 60
estudantes ocuparam salas do CB por mais de dois anos, reivindicando uma moradia
gratuita aos estudantes de baixa-renda. Essa ocupação estudantil ficou conhecida como
"A Taba", em referência a tabas indígenas que são caracterizadas por uma cultura de
coletividade.

O reitor da época, Paulo Renato Costa Souza, prometeu que 1.500 vagas –
correspondente a cerca de 10\% do total de estudantes daquela época – seriam
entregues até o mês de julho de 1989 e, no caso de não cumprimento das cláusulas que
garantiam a construção da moradia, ele assinou um documento que autorizava o DCE
Unicamp a ocupar novamente o campus, assim como o Movimento TABA havia feito. Em
1990, a Moradia foi inaugurada com novecentas e quatro vagas, já abaixo do prometido.

Na greve de 2016, os estudantes reivindicaram a expansão das vagas e conseguiram
que o reitor José Tadeu Jorge assinasse um documento que a garantisse. Com o fim da
greve, criaram-se grupos de trabalho (GTs) para analisar e avaliar a ampliação e os
programas de permanência estudantil da universidade. Eventualmente, os GTs foram
dissolvidos e o terreno, que estava em processo de compra, ficou embargado em
cartórios. Finalmente, membros da administração central da universidade admitiram em
reunião com membros do movimento estudantil que a nova moradia "não vai sair".

Além disso, há diversos problemas na infraestrutura que indicam uma falta de
manutenção generalizada, como por exemplo estruturas que foram evacuadas sob
perigo de desmoronamento; problemas elétricos, de drenagem nas ruas e de infiltração
de água nas casas; dificuldades na separação do lixo reciclável; e vazamentos. Há
muitos casos de superlotação, com casas abrigando até sete estudantes, e vários casos
de furto, principalmente de bicicletas dos moradores.

Uma moradia estudantil digna é um espaço determinante para a garantia da
democracia no acesso e permanência no Ensino Superior - incluindo os campi que
sequer têm programa de moradia estudantil, em Piracicaba e em Limeira.

Por fim, há uma disputa de projeto sobre a moradia: ao invés de construir mais casas,
a reitoria tem "investido" na Bolsa Moradia - um valor de cerca de R\$ 400 pago a
estudantes para se manterem através de contratos de aluguel. O impacto na organização
estudantil é enorme, pois es bolsistas estão pulverizades e não têm um espaço para se
encontrar e auto-organizar, caso precisem defender suas pautas.

Os milhões de reais por ano vão pelos ares em aluguéis que em pouco tempo poderiam arcar com a construção de uma moradia inteira. Crises financeiras cada vez mais constantes - lembre-se que ainda estamos vivendo consequências da "crise de 2008" -
podem pôr em xeque essas bolsas, o que não ocorreria numa moradia construída. O tiro
cego da reitoria acerta, contudo, o mesmo alvo: a especulação imobiliária de Barão, que
só tem a agradecer.

\vspace{1em}
\textbf{O que você não viu:}
\vspace{-0.5em}
\begin{itemize}
    \setlength\itemsep{-0.2em}
    \item \textbf{1986}: Sem ter onde morar, estudantes estabelecem o movimento TABA e acampam no CB até a conclusão da construção da moradia.
    \item \textbf{2016}: Greve estudantil é lançada e dura 5 meses, ocupando a reitoria.
    \item \textbf{2017}: Longa sessão do CONSU aprova cotas étnico-raciais e sociais, diante de um ato com mais de 600 pessoas.
    \item \textbf{2019}: Es primeires estudantes cotistas entram na universidade. Estudantes organizam uma rede de apoio solidário para viabilizar que indígenas e cotistas tenham condições suficientes para se estabelecer nas primeiras semanas, o que ocorreu também em 2020.
    \item \textbf{2023}: Conquista de expansão...
\end{itemize}

<imagens>

\section{SAE - Serviço de Apoio ao Estudante}

O SAE, principal órgão de apoio a estudantes na Unicamp, atua em várias frentes de
assistência estudantil. Esta se dá por meio do gerenciamento de bolsas-auxílio,
assistência social e orientações educacional/jurídica/psicológica, além do apoio a
projetos acadêmicos e sociais e programa de intercâmbio de estudantes no exterior.

O SAE também é responsável pela gestão de estágios na Universidade.

\textbf{Localização}: prédio do Ciclo Básico, 3o piso (responsável pelo gerenciamento de
convênios, estágios, bolsa-pesquisa e bolsa-empresa); e ao lado da cantina do DCE
(serviço social, responsável pelo gerenciamento de bolsas-auxílio).

\textbf{Horário}: de segunda a sexta, das 8h30 às 20h, no período letivo.

\begin{tags}
    \site{www.sae.unicamp.br} \sep \email{sae@unicamp.br}
\end{tags}

\section{Bolsas-Auxílio e Inscrição}



\section{Orientação Jurídica Gratuita}

O \textbf{SAE, Serviço de Apoio ao Estudante}, oferece orientação jurídica gratuita para
uma porção de casos, não necessariamente relacionados à universidade ou vida
acadêmica: o serviço está aí justamente para orientar estudantes nas questões pessoais
pelas quais a gente pode passar.

\textbf{Direito Civil}: contratos em geral, contratos de locação, acidentes de trânsito, ação
revisional de aluguel, divórcio, pensão, etc.

\textbf{Direito Penal}: violência contra a pessoa, lesões corporais, furto, roubo, etc.

\textbf{Direito do Trabalho}: caracterização de relação de emprego para não registrades e
direitos trabalhistas em geral (CLT).

Procedimento: e alune deve se dirigir pessoalmente ao SAE \textit{(Rua Sérgio Buarque
de Holanda, 251 - Cidade Universitária, Campinas - SP) para obter os esclarecimentos
desejados. Delitos de consumo: encaminhar ao SEDECON ou ao PROCON.}

\section{Atendimento médico e odontológico - CECOM}

O CECOM é o Centro de Saúde da Comunidade e oferece atendimento ambulatorial
gratuito aos seus alunos em diversas especialidades médicas, saúde mental, fisioterapia,
nutrição, enfermagem e odontologia..

A CSS, Coordenadoria de Serviços Sociais, é responsável pelo planejamento e
execução de programas de saúde voltados à comunidade universitária da Unicamp –
alunes, funcionáries e docentes.

Na prática é quase como um plano de saúde da Unicamp. Você pode marcar
consultas médicas e fazer exames. O CECOM fica próximo ao HC (Hospital das Clínicas)
e para chega lá você pode pegar o circular interno, é o ponto final de algumas linhas.

Fique de olho na seu e-mail institucional perto do período de inverno, pois \textbf{o CECOM
costuma disponibilizar vacina contra gripe gratuitamente} - uma dose em clínica
particular é bem cara. O CECOM também oferece vacinas contra COVID, que são
obrigatórias e se você não estiver com elas em dia pode ter problemas para se
matricular.

Outra coisa que você não pode perder no CECOM são os testes de ISTs (infecções
sexualmente transmissíveis). O teste cobre várias ISTs, inclusive o HIV, e o resultado sai
na hora. É necessário agendar, exceto se você tiver tido exposição a alguma IST.

Para marcar consultas com dentista, vá ao CECOM e peça informações. Basicamente é
embaixo do CECOM, muito fácil de entender como chegar se alguém apontar com e
dedo e disser "ali". O atendimento funciona muito bem e já começa com uma palestra
instruindo como cuidar bem da boca no dia a dia. Afinal, a melhor saúde é a preventiva.
Mas se você estiver com dores, te atendem na hora sem marcar consulta nem assistir
palestra.

Veja todas as informações sobre atendimento e contato no site do CECOM:
\url{www.cecom.unicamp.br}.

\section{Hemocentro}

Essa é para quem é ou para quem quer ser doadora ou doador de sangue. O centro de
hematologia e hemoterapia (Hemocentro) é o órgão da Unicamp responsável pela coleta
e doação de sangue.

Para ser doadore voluntárie, é necessário portar um documento com foto e seguir os critérios para a doação de sangue. É realizada uma entrevista que tem como objetivo
caracterizar riscos para e receptore e verificar contra-indicações a doação. Para saber
mais sobre o processo, é só visitar o site: \url{www.hemocentro.unicamp.br}.

O setor de Hemoterapia desenvolve campanhas permanentes de conscientização
sobre a doação voluntária de sangue em municípios da região, visando o engajamento
de toda a comunidade em seu trabalho e na manutenção de estoques seguros para o
atendimento das necessidades transfusionais dos pacientes.

O Hemocentro fica localizado acima do HC, próximo ao CECOM. Portanto, se quiser se
deslocar até lá, usar o circular interno é uma boa opção.

<imagens>

\section{Cuidando da Saúde Mental}

Vamos ser sinceres contigo: na faculdade, a gente tem muita coisa boa e muita coisa
ruim também. Pronto, falamos. Nada é perfeito, né? Pela relevância que o curso toma nas
nossas vidas, isso tem um peso que pode acarretar vários efeitos psicológicos na gente.

Levantamentos demonstram altos índices de estresse, ansiedade e depressão na pós-
graduação. Calma, você está lendo o manual certo: é que muitos dos fatores que tornam

o ambiente acadêmico insalubre na pós também estão presentes na graduação.
Por isso, o melhor que temos a fazer é entender isso e saber o que fazer pra encarar
esse problema sem sofrer demais com ele.

\subsection{Estudar dá (ou é) trabalho}

Ao contrário do que dizem, universidade não é só rolê. Primeiro porque a gente
estuda, e MUITO. É uma carga de cobrança imensa e, honestamente, injusta. Infelizmente,
a maioria de nosses professores nas Exatas não têm formação em licenciatura ou
pedagogia – o que es deixam sem importantes ferramentas de didática. Aí, alguns
passam muita matéria e acham que, se pouca gente reprovou (você leu certo, reprovou),
é porque alguma coisa está errada. Os catálogos dos cursos também têm problemas.
Assim, o risco é grande de a gente acabar com a sensação de falta de tempo,
ansiedade, dificuldade de nos relacionarmos e até mesmo sentimento de culpa quando
descansamos ou nos divertimos (aquele "nossa, eu devia estar estudando"). Parece que
somos incapazes por causa do alto nível das matérias, mas na verdade existe sobretudo
um baixo nível do ensino, em várias áreas e institutos. Para quem já está começando a
vida adulta, tem que ainda pagar contas e lidar com a distância da família, isso também
acaba exercendo bastante pressão. Pesa ainda o quanto a nossa cultura negligencia a
saúde mental, até porque o acesso universal a psicólogues infelizmente não está
disponível no sistema de saúde gratuito e, ainda que tentemos optar por pagar, eles são
caros.

\subsection{You Are (Not) Alone}

É no espaço da universidade que estamos submetidos a essa infinidade de cobranças,
a conflitos para os quais não estamos preparados e, como em todo espaço em que
temos que enfrentar a hostilidade do mundo, você pode se sentir só. No entanto, é
também nesse espaço que podemos fazer amizades, encontrar companheires que nos
ajudam nos momentos mais difíceis. Então, não se esqueça que a faculdade não é e
jamais deveria ser uma competição. Você tem como ajudar seus colegas e tem onde
buscar apoio quando precisar, tem colegas e veteranes que passaram ou passam pelas
mesmas coisas que você. Aja com franqueza e sinceridade para consigo mesmo, busque
entender o que acontece com você e se comunique, sem medo de pedir ajuda ou
desabafar. Ter dificuldades não te faz melhor ou pior e o \textbf{seu rendimento acadêmico não
significa tudo na vida}. Nosso sistema de educação é terrivelmente falho e de forma
alguma consegue dar a última palavra sobre o que uma pessoa é ou deixa de ser. Em
suma, bitolar não é saudável!

Por isso, não cobre muito de si e não se negligencie. Busque o que te fortalece, conte
com a comunidade, converse com o CACo e as demais entidades, apresente ideias, se
abra ou mesmo denuncie aquilo que te faz mal. Juntes podemos melhorar a nós e nosso
entorno. Se for preciso, lutaremos para que melhore! Nesse mundo que individualiza
tudo, lembre-se que a saúde mental é uma questão coletiva.

\subsection{SAPPE - Serviço de Apoio Psicológico e Psiquiátrico}

A universidade oferece serviços de assistência psicológica e psiquiátrica gratuitamente através do SAPPE, órgão ligado à Pró-Reitoria de Graduação (PRG). Ele oferece
algumas modalidades de atendimento, divididas em duas principais categorias: o
atendimento regular e o pronto atendimento.

O atendimento regular é um tratamento mais contínuo que é feito geralmente em
quatro sessões, a menos que precise ser aumentado, de acordo com a opinião
profissional. Para utilizar o serviço você deve preencher uma Ficha de Cadastro para
atendimento regular de acordo com os horários disponíveis no site do SAPPE (disponível
abaixo) na recepção do Serviço (endereço e instruções de chegada abaixo) e agendar a
participação no Grupo de Recepção. Após a participação no grupo, você pode agendar
uma Entrevista de Triagem que te encaminhará para o tratamento adequado.

\textbf{Observação}: o Grupo de Recepção é uma palestra que explica o funcionamento do
serviço, mas ela parece ser montada para te desmotivar a usá-lo. Os períodos de espera
entre cada etapa também podem ser bem longos e servem para esse mesmo fim. Então,
se você realmente sente a necessidade de buscar esta orientação, \textbf{não desista}.

O SAPPE também tem pronto atendimento. Essa é uma sessão única voltada para
quando você estiver passando por uma crise ou emergência, mas com os horários
limitadíssimos. Nesse caso, você deve passar na sede do SAPPE nos dias e horários
disponíveis no site.

\textbf{Localização}: R. Sérgio B. de Holanda, 251, 1o piso Em frente ao prédio da DAC; mesmo
prédio do SAE e CLE, ao lado do Bandejão.

%\textbf{Horário}: de segunda a sexta, das 8h30 às 20h, no período letivo.

\begin{tags}
    \tel{19 3521-6643} ~~~ou~~~ \tel{19 3521-6644} \\
    \facebook{SAPPE-Unicamp-112906617163477} \sep \email{sappeass@unicamp.br}
\end{tags}

<imagem>

\section{Transporte interno}

Com uma infraestrutura do porte de nossa universidade, há muitos recursos dedi-
cados a uma rede de transporte interno dentro e fora do campus de Barão Geraldo, de
graça!

O Circular Interno funciona dentro do nosso campus e tem várias rotas cobrindo os
lugares mais distantes, inclusive o IC. A partir das 18h até as 23h funciona somente o
Circular Noturno, que tem a frequência bem menor e menos rotas.

Também é livre o acesso de qualquer estudante ao ônibus da Moradia, comumente
utilizado por estudantes que moram lá ou em repúblicas nos arredores da Av. Santa
Isabel. Basta aguardar no ponto de ônibus no bolsão de estacionamentos da Biblioteca
Central, ao lado do Bandejão, e mostrar o RA. De lá também sai o Intercampi, ônibus que
conecta os campi de Barão e de Limeira. O uso é limitado, e você tem que se inscrever
com antecedência no site do serviço.

\textbf{Você pode conferir as rotas, horários e até posição geográfica ao vivo pelo aplicativo
Unicamp Serviços, que também tem várias outras coisas legais.}

Além desses transportes oferecidos pela universidade, quem é de São Paulo também
tem a opção de utilizar o Massa Crítica, que é particular. Ele existe para suprir a
demanda de transporte entre a Unicamp e a USP e seus arredores sem passar pela
rodoviária de Campinas e da capital.

\section{O que comer na Unicamp}

\subsection{Histórico do Bandejão}

A universidade, após sua fundação, passou 10 anos sem ter bandejão. O primeiro foi o
RU (Restaurante Universitário); em seguida, um prédio que pertencia ao CABS (Centro
Acadêmico Bernardo Sayão, da Engenharia Elétrica) foi convertido no RA (Restaurante
Administrativo), perto da FEEC; depois foi construído o RS (Restaurante Saturnino), mais
moderno e perto do IC. Hoje todos os campi da universidade têm acesso a bandejões,
inclusive os colégios técnicos (COTUCA em Campinas e COTIL em Limeira).

O preço das refeições é subsidiado e, sem ele, cada estudante teria que viver na base
da marmita, passando maus bocados para ter tempo para cozinhar ou pagar o preço
comercial de todas as refeições.

O Bandejão começou a servir café da manhã a pouco tempo e é uma reivindicação
quase tão antiga quanto a abertura nos finais de semana - uma pauta importantíssima
para quem não pode voltar pra casa da família nos fins de semana e tem que se virar
por aqui, mas que ainda não conseguimos conquistar.

\subsection{Sobre o Bandejão em si – ou Bandeco, Bandex, RU, RA, RS, etc!}

Um dos momentos de glória do dia de ume future engenheire, cientista ou bacharel é
o Bandejão. É a hora de intensas e indiscutíveis emoções. Jamais tente descobrir o sabor
do suco pelo paladar, é mais fácil ler no cardápio do dia. O Bandejão é o lugar de você se
encontrar com amigues (combinando antes ou não), contar os micos nas aulas e jogar
conversa fora.

Eles funcionam de segunda à sexta e são bem afastados um do outro, o que torna
qualquer instituto razoavelmente perto de um bandeco, independente de onde tenha
sido ou vai ser sua aula. Em períodos especiais, como férias, os restaurantes podem
funcionar em horários reduzidos ou não abrirem, então fique de olho no site da
prefeitura, no GDE, ou principalmente no aplicativo da Unicamp. Nessas plataformas
você também pode encontrar o cardápio da semana, que conta com a famigerada carne
picada pelo menos um dia.

Atualmente, o preço do café da manhã é de R\$ 2,00 e do almoço e jantar é de R\$
3,00. Sem dúvida ainda é um ótimo custo benefício, porém ele já foi melhor:
historicamente, o preço do bandejão era de R\$ 2,00. No entanto, entre 2017 e 2018,
numa reunião extraordinária do Conselho Universitário, foi aprovado o aumento das
refeições dos restaurantes universitários. Por mais que a inflação tenha aumentado, as
bolsas estudantis não tiveram um reajuste equivalente.

Existem outras opções para se alimentar dentro da Unicamp. Contudo, elas
certamente custam mais do que R\$ 3,00.

\subsection{Cantinas}

As cantinas oferecem algo mais voltado para o "lanche da tarde", ao invés de
refeições completas, como salgados e bolos. No entanto, existem algumas com pratos
feitos ou com self-service por quilo, saindo bem mais caro que o bandejão. Elas eram
bem numerosas no passado, porém conforme suas licenças de funcionamento foram
expiran-do, a universidade não abriu novos editais para renová-las e ainda não se
pronunciou a respeito. Por conta disso, nem todos os institutos tem cantinas e não
sabemos ao certo quantas cantinas estarão funcionando nesse ano. O mesmo fenômeno
ocorre com as barraquinhas de impressão e xerox dos institutos.

\subsection{Feirinha}

Carinhosamente apelidada de feirinha, é uma feira que ocorre na praça do Ciclo
Básico às quartas e quintas (quando temos feriados, os dias mudam). A variedade de
opções é imensa: existem barracas de suco, tapioca, macarrão, açaí, pastel, doces,
strogonoff, entre muitas outras – além de produtores de artesanato. Também há
restaurantes e supermercados em Barão, veja na seção Vivendo em Barão.

\section{Mantenha-se em segurança}

\subsection{Campus Tranquilo}

O Campus Tranquilo era um programa de segurança do campus implementado após
formulação com a comunidade acadêmica e profissionais de segurança pública como
alternativa, em 2013, à tentativa da Polícia Militar de estabelecer presença permanente
no campus. O argumento é que a polícia pode cumprir um papel truculento e repressivo
- sobretudo em protestos contra cortes de verbas. Afinal, mais de uma vez, em várias
universidades, como na USP, funcionários e estudantes tomaram bomba de gás e foram
detidos pela PM diante de votações importantes do Conselho Universitário.

Após uma decisão judicial, funcionários da FUNCAMP treinados no programa foram
demitidos em massa e a universidade re-terceirizou o serviço. Essa força de trabalho
agora têm condições mais precárias e muito menos treinamento. O resultado é que a
guarda patrimonial se tornou mais policialesca, com casos de abordagens agressivas.

O patrulhamento patrimonial continua mas a maioria des guardinhas se apresenta
amigavelmente, sendo possível pedir escolta durante a noite. Também há uma
ambulância em caso de emergências de saúde.

\subsection{Botão de pânico}

O Botão do Pânico é um app disponibilizado pela própria Unicamp, que tem como
objetivo registrar situações de pânico dentro do campus. O app está disponível tanto na
Play Store quanto na App Store, sendo que os credenciais para login são seu RA e a sua
senha da DAC. O ideal é já fazer o download do aplicativo e fazer o login, pois numa
situação de risco não haverá tempo para isso. O aplicativo utiliza sua localização para
pedir emergencialmente a presença de algum responsável da segurança da Unicamp.

\subsection{Grupo da Unicamp no Facebook}

O grupo da Unicamp no Facebook tem vários objetivos, como integração e diversão,
no entanto também pode ser utilizado para divulgação de informações. De vez em
quando rolam no grupo algumas publicações sobre uma situação de risco que ocorreu
em Barão Geraldo, explicando o que aconteceu e onde especificamente aconteceu. Por
isso, fique de olho!

\subsection{Página Avisa as Minas}

As dependências de Barão Geraldo também são perigosas para mulheres. Para não só
evitar, mas também combater situações de assédios verbal e físico, perseguição e
ameaças, foi criada a página "Avisa as Minas de Barão/Unicamp" (tagueado por \#AMBU),
para deixar todo e qualquer tipo de aviso que você ache importante sobre essas
questões de segurança citadas. Nela, é possível postar alertas o mais rápido possível,
então lembre sempre de ativar para receber as notificações e mandar nos grupos de sala
para que cada vez mais meninas sejam alertadas.

\chapter{Estudando na Unicamp}

\section{Siglas!}

Antes de qualquer coisa, você vai logo perceber que, dentro da Unicamp, existem um zilhão de siglas diferentes- Não se preocupe em decorar tudo!!! Com o tempo, você vai se familiarizar bem com elas, prometo!

Separamos, nessa sessão, explicações sobre algumas siglas, códigos e outros termos que são muito usados dentro da Unicamp- Assim, enquanto estiver lendo o manual ou vivendo sua vida de ingressante, terá em mãos um glossário para consultar! Se houverem mais dúvidas ou siglas que não constam aqui, não hesite em perguntar a sues veterenes ou a nós do Centro Acadêmico. ;)

\begin{itemize}
    \item \textbf{Currículo pleno}: é o conjunto de disciplinas do curso que e alune tem que cursar para se formar.
    \item \textbf{Créditos}: unidade elementar de horas-aula de qualquer curso da Unicamp. Um crédito equivale a uma hora-aula semanal, ou a 15 horas-aula semestrais. Cada disciplina possui um “número de créditos” equivalente a ela- Normalmente, esse número de créditos indica quantas horas semanais você terá de aulas dessa disciplina- Por exemplo: cálculo 1 é uma matéria de 6 créditos; você terá, portanto, 3 aulas de 2 horas por semana (ou 6 horas-aula por semana). A cada semestre, você irá acumulando créditos. Para se formar no curso, você precisa completar um número específico de créditos (explicitados no currículo pleno).
    \item \textbf{CR (Coeficiente de Rendimento)}: valor entre 0 e 1 da média das notas em todas as disciplinas cursadas ponderada pelos créditos (ou seja, ir mal numa matéria de 6 créditos pode prejudicar muito seu CR). PS: atenção redobrada se sua intenção é seguir carreira acadêmica e/ou intercâmbio (mais detalhes sobre isso à frente).
    \item \textbf{CP (Coeficiente de Progressão)}: parte do curso que você já cumpriu (CP = 0,6123 significa que se cumpriu 61,23\% do curso).
    \item \textbf{CPF (Coeficiente de Progressão Futuro)}: é o CP que você terá no fim do semestre caso passe em todas as disciplinas que você se matriculou.
    \item \textbf{CPE (Coeficiente de Progressão Exigido)}: usado para fins de cancelamento de matrícula. Para que o aluno possa continuar a fazer o curso, ele precisa ter um CP maior ou igual ao CPE daquele semestre.
    \item \textbf{GDE}: incrível facilitador de vidas! É um site, e também rede social, onde você pode usar seu login da DAC para navegar entre disciplinas, salas de aula e listas de colegas. O Planejador e a avaliação de professores (extra-oficial) são a principal ferramenta. Projeto de Felipe Guaycuru, um ex-aluno da Engenharia de Computação, reconhecido pela DAC. O link é \textbf{\url{gde.ir}}.
    \item \textbf{Pré-requisito}: matéria(s) que precisa(m) ter sido cursada(s) para que se possa fazer outra(s) matéria(s). Existem dois tipos de pré-requisitos: Os pré-requisitos totais, mais comuns, do qual é exigido tanto a aprovação por nota como por frequência e os pré-requisitos parciais, mais raros, do qual o aluno não precisa ter sido aprovado por nota, mas tem que ter tido aprovação por frequência e nota final maior ou igual a 3,0. Os pré-requisitos parciais são identificados com um asterisco na frente do código da disciplina.
    \item \textbf{AA4xy}: um tipo de pré-requisito. Não se trata de nenhuma disciplina. Para fazer disciplinas com esse pré-requisito, e alune tem que tem um CP maior ou igual a 0,xy.
    \item \textbf{AA200}: outro tipo de pré-requisito existente, mais presente em disciplinas eletivas. Também não se trata de nenhuma disciplina. É apenas uma autorização da coordenadoria do curso. Se sobrar vagas para a disciplina e a coordenadoria do curso for com a sua cara, você faz a disciplina.
    \item \textbf{CCG (Comissão Central de Graduação)}: órgão colegiado da Unicamp, é encarregada da orientação, supervisão e revisão periódica do ensino na Universidade. Cabe recurso à CCG de quaisquer decisões das Unidades afetando o ensino.
\end{itemize}

\section{Disciplinas}

\subsection{Matrícula}

Com exceção do primeiro semestre letivo, no qual você já entra matriculade em todas as matérias obrigatórias, \textbf{na Unicamp é você quem vai ter que montar sua grade de disciplinas}. A cada semestre, você deve enviar à DAC um requerimento de matrícula com as disciplinas que deseja cursar no semestre. Mas não é nada complicado, e ninguém vai te deixar esquecer! Você pode consultar as matérias sugeridas por semestre no catálogo de curso, disponibilizado pela DAC, na sua “Proposta para Cumprimento de Currículo”.

Para te auxiliar na parte mais legal do semestre – o planejamento da grade –, existe o GDE (\url{gde.ir}): uma ferramenta criada por um aluno da engenharia e adotada pela DAC que disponibiliza um planejador já com as matérias sugeridas para o semestre e as turmas. Também é uma rede social interna, permitindo avaliação de oferecimentos de matérias e de professores, que serve de feedback às outras pessoas.

Recomendamos que peça ajuda a veteranes quando for montar seu horário! Informe-se sobre todes es professores que oferecem as matérias, se são mais tranquiles ou pegam no pé, se dão aula bem ou mal, se demoram para entregar as notas... Você vai poupar muita dor de cabeça! Os melhores lugares para essas discussões são os grupos de Facebook da Unicamp, grupos com veteranes no wpp/telegram e o servidor de \textit{discord Corredores do IC} – por onde centramos nossa comunicação durante o EAD. Dá um alô pra gente que te mandamos o link para entrar!

Porém, nem sempre você vai conseguir exatamente o que quer na matrícula. É possível que te joguem para uma turma em outro dia da semana, caso a turma que você pediu lote. Isso é muito comum e, nesses casos, pode tentar de novo no \textbf{período de alteração de matrícula}, que acontece, normalmente, próximo ao início das aulas. Nesse período, é possível pegar uma matéria que não conseguiu na matrícula normal, permutar a turma de uma matéria que já se matriculou ou até mesmo desistir de alguma delas, sem nenhum tipo de penalidade.

\subsection{Datas e prazos}

Para não ser pegue pelo arrependimento de ir a alguma aula em dia de folga ou de não poder mais trancar uma matéria que considera já perdida, anote as datas importantes todo início de semestre no seu calendário. Você pode encontrar essas datas no calendário da DAC: o site possui um filtro que seleciona o que é de interesse de alunes de graduação, da pós, de docentes e de funcionáries – o que torna sua vida muito mais fácil.

Isso não serve apenas para saber quais são os feriados: serve também para períodos de estudo para o exame – assim como o período para a aplicação da prova –, período para desistência de disciplinas, trancamento de matrícula, férias e matrícula no próximo semestre. Mesmo seguindo o calendário, às vezes há mudanças no calendário no meio do período! Fique atente!

\section{Cancelamento, trancamento e desistência}

Embora praticamente todes estudantes da Unicamp usem esses três termos
indiscriminadamente, como se fossem sinônimos, eles têm significados bastante
distintos para a DAC. Aí vai:

\subsection{Desistência de matrícula em disciplinas}

A desistência é o processo que é chamado peles alunes de “trancamento”. Você deixa de cursar essa disciplina no meio do semestre, tendo de cursá-la em algum semestre posterior (se for obrigatória). Só é possível desistir uma vez da disciplina e pode-se pedir desistência até que se tenha passado metade do semestre (confira a data limite da desistência no calendário!).

\subsection{Trancamento de matrícula}

Processo em que você não cursa nenhuma disciplina da Unicamp durante o semestre. É possível fazer até dois trancamentos de matrícula (em semestres seguidos ou não) e não se pode trancar nenhum dos dois semestres do ano de ingresso. A desistência de todas as disciplinas configura-se como trancamento. O trancamento é pedido na DAC, e pode ser pedido até que se tenha transcorrido 2/3 do semestre (de novo, veja o calendário!). Para cada trancamento, o prazo máximo de integralização é postergado, ou seja, o seu tempo fora da Unicamp não contará para o progresso do curso – isso pode ser muito útil em situações críticas.

\subsection{Cancelamento de matrícula}

Processo em que você se desliga da Unicamp, seja por motivo de jubilamento; por ter faltado às duas primeiras semanas do ano de ingresso; por ter sido reprovade em todas as disciplinas do primeiro ou do segundo semestre de ingresso; por ter sido expulse; por cursar outra universidade pública simultaneamente ou por vontade própria.

\subsection{Jubilamento? Mas eu acabei de entrar!}

É importante garantir a sua permanência na vaga! Segundo o regimento, o ingressante que tiver ausência em todas as aulas nas duas primeiras semanas do primeiro semestre terá a matrícula cancelada. Além disso, a não aprovação em nenhuma


em todas as aulas nas duas primeiras semanas do
primeiro semestre terá a matrícula cancelada. Além disso, a não aprovação em nenhuma

matéria obrigatória do curso em algum dos dois primeiros semestres resulta em jubilamento, independente do número de matérias eletivas ou completadas em outro semestre. Mantenha um CP (Coeficiente de Progressão) razoável: um CP menor que o CPE (Coeficiente de Progressão Exigido), também jubila (calma! tem que ser menos de metade dos créditos previstos). De resto, é sempre bom dar uma olhada nas outras regras no site da DAC, mas nada pra se preocupar demais agora no começo do curso. Elas podem ser encontradas no artigo 49 do regimento de graduação, disponível no site da DAC: \url{www.dac.unicamp.br/portal/graduacao/regimento-geral}.

\subsection{Mudança de catálogo}

Para cada ano, existe um catálogo correspondente com todas as disciplinas que os alunos devem fazer para se formar (currículo pleno). O seu catálogo é o 2023, por conta do seu ano de ingresso, mas é possível alterá-lo. Como a computação é uma área muito dinâmica, ocorrem várias mudanças no decorrer dos anos, como adição e remoção de disciplinas e aumento ou diminuição de créditos.

Como essas mudanças ocorrem num catálogo específico de um ano e só valem a partir desse, elas não são aplicadas a estudantes de anos anteriores. Para mudar isso, é possível mudar o ano de catálogo. Porém, existem ressalvas: se você mudar de catálogo, pode ser que você precise fazer mais matérias para se formar. Há também o risco de não conseguir equivalência de algumas matérias que são semelhantes entre si nos catálogos. Para mudar de catálogo, é só preencher o formulário \url{bit.ly/1LHPdyG} e entregar à DAC. Mas confira sempre com seu Centro Acadêmico e com veteranes se vale a pena!

\textbf{> Como ocorrem essas mudanças de catálogo?}

Todas feitas com muita discussão entre professores e representantes de alunes. Se
você não estiver antenade, pode ficar sabendo nas reuniões semestrais de avaliação de
curso. Fazemos nosso melhor para informar sobre todas as alterações e colher sugestões
nós mesmes, então fique atente!

\subsection{Avaliação de curso}

Avaliação de curso é uma reunião anual proposta pela coordenadoria que visa discutir todos esses pontos. Participe, pois é sua chance de tornar o curso melhor falando direta ou indiretamente com es coordenadores! Caso tenha sugestões, converse com a gente também e faremos o necessário para defender melhorias no curso!

Falamos um pouquinho mais sobre isso adiante, em \textbf{Avaliação de Professores}.

\section{Eletivas, proficiência e aproveitamento}

A Unicamp oferece a oportunidade de personalizar seu currículo de acordo com seu interesse por meio das disciplinas eletivas. Ao contrário das disciplinas obrigatórias, com as eletivas você pode escolher a matéria que vai cursar. Alguns créditos podem ser cumpridos com qualquer disciplina oferecida pela Universidade, outros estão restritos a um determinado conjunto. Para mais detalhes, consulte seu catálogo. Você pode socializar com outros cursos, aprender uma língua estrangeira, assistir a seminários ou obter um certificado de estudos na FEEC ou no IC.

\textbf{Teste de proficiência} é uma prova que permite dispensa de cursar uma disciplina (desde que você obtenha a nota mínima, é claro). Se você acha que sabe o suficiente sobre eletromagnetismo, por exemplo, pode tentar a proficiência de Física Geral III. Nem todas as disciplinas oferecem o teste e você só pode fazê-lo uma vez por disciplina – e se você já se matriculou na disciplina e não passou, não pode fazer. Para as matérias específicas de computação, a proficiência é mais voltada para quem já estudou ou trabalhou com programação, pois requer alguma experiência. Além disso, fazer o teste de proficiência também é obrigatório para se matricular nas disciplinas de língua inglesa e japonesa, independentemente de conhecimento prévio na língua.

Fique ligade no calendário da DAC para não perder as datas de inscrição! As datas dos testes de línguas são sempre no começo do ano, diferentes das demais, que são no fim de cada semestre. 

Disciplinas eletivas e teste de proficiência estão relacionados porque muitas pessoas fazem proficiência em disciplinas de línguas – eliminando créditos de eletivas, em alguns casos para evitar o jubilamento, em outros para não ter que passar mais um semestre na faculdade. Converse com veteranes para descobrir o melhor jeito de usufruir dessa liberdade que poucas universidades oferecem! Dificilmente você não encontrará algo com o qual se identifica ou que não ensine lições interessantes.

Por fim, caso você tenha iniciado um curso de graduação antes e queira aproveitar os créditos no novo curso, existe o \textbf{aproveitamento de estudos}. Se as disciplinas foram feitas na UNICAMP esse processo é automático. Se foi em outra instituição, a partir do seu pedido, é feita análise e comparação das ementas das disciplinas. Atenção: só pode ser feita no 1o semestre!

Para mais informações sobre proficiência e aproveitamento, acesse: \url{bit.ly/2BS5O6e}.

\section{Centro de Ensino de Línguas - CEL}

O CEL, como o nome já diz, é o órgão responsável por oferecer aulas de diferentes  idiomas a alunes da Unicamp. Seja como disciplinas obrigatórias (para a engenharia de  computação, Inglês Instrumental I é uma delas) ou como eletivas. O CEL possui turmas  de inglês, espanhol, alemão, francês, japonês, dentre outras línguas.

Todas as línguas no CEL oferecem níveis diferentes; assim, se você souber um pouco  de Francês, não precisa começar do início. É só realizar o teste de proficiência, oferecido  no primeiro semestre de cada ano, e você pode avançar algumas turmas (ou todas) em  uma determinada língua. Também é possível fazer o teste de proficiência pra eliminar  disciplinas (o que a maioria das pessoas acaba fazendo com Inglês Instrumental). Todes  es alunes que quiserem se inscrever em \textbf{Inglês I (LA112)} ou \textbf{Japonês I (LA111)} devem  fazer o teste, mesmo que não saibam nada da língua; mas ele não é necessário para  matrículas nas outras línguas. 

A matrícula em disciplinas do CEL é feita de maneira natural, na DAC, durante o  período de matrícula. 


\section{Avaliação de professores}

Achou que e professore ensinou muito mal? Falou da vida, do universo e tudo mais  ,exceto da disciplina? Foi incoerente? Ou, pelo contrário, achou e professore incrível? Se você sentir que as coisas simplesmente não estão funcionando em sala de aula ou  que algo de errado não está certo, traga o assunto para o CACo. Seu problema com a  matéria, na maioria das vezes, não é só seu. Como representantes dos seus interesses,  temos o dever e o prazer de fazer tudo que estiver ao nosso alcance para melhorar a  situação. 

O GDE também tem um sistema de avaliação de professores, cuja nota costuma ser  usada peles alunes como um dos critérios no momento de decidir com que professore  puxar uma matéria – mas prefira o feedback de veteranes. 

Além disso, temos a \textbf{Reunião de Avaliação de Curso}, já comentada. A data e o horário  são divulgados pelas unidades e pelo CACo. Essa é uma oportunidade de passar para às  coordenadorias do curso não só suas impressões sobre professores e disciplinas, mas  sobre qualquer assunto relacionado ao curso. Complementarmente, todes professores  devem disponibilizar um formulário de avaliação nas últimas aulas de cada semestre.  Esse é o momento para que você possa separar os acertos dos erros, portanto preencha  com seriedade. Os dados serão analisados pelas Comissões de Graduação de cada  unidade e os comentários escritos serão repassados para e professore. 

Tenha sempre em mente que a nossa percepção sobre o oferecimento de uma  disciplina não é óbvia para es professores. Preencha todas as avaliações com  sinceridade, use sempre os espaços dedicados a comentários com críticas construtivas. 


\section{CR e Reprovação}

Durante o curso, você provavelmente vai ouvir que se preocupar com seu CR é  bobagem; que estudar para tirar nota não leva a lugar nenhum; que depois de formade  não é seu CR que te colocará no mercado de trabalho; etc. \textbf{Cuidado! Trabalhar numa  empresa não é a sua única opção de vida após formade.} E atenção, pois “após formade”  não significa durante o curso. 

Durante o curso, você vai ter a possibilidade de participar de várias atividades  acadêmicas e algumas delas vão exigir bom aproveitamento acadêmico. \textbf{Por exemplo, para se candidatar à monitoria de uma disciplina, você precisa ter CR acima do da média  da turma. Para pleitear uma bolsa de iniciação científica, onde há concorrência entre alunes do país todo, também será exigido bom aproveitamento – assim como para uma  bolsa de mestrado.} Caso você não saiba, mestrado faz parte da pós-graduação, ou seja, o  seu CR vai te influenciar até após formade. 

Cuidado também com a reprovação. Há instituições, como a \textbf{FAPESP} (Fundação de  Amparo à Pesquisa do Estado de São Paulo – que é a maior fomentadora de pesquisas  do estado de SP e que paga os maiores valores de bolsas do país) \textbf{que te excluem de  qualquer disputa só por ter uma reprovação no seu histórico escolar da graduação}. Não  te excluem oficialmente, mas como é muito concorrido por ser a melhor pagadora, seu  nome vai para o final da lista. 

Parece óbvio para quem estuda tira boas notas, mas até você aprender a estudar  como a universidade exige, pode demorar um pouco; e há pessoas que nunca aprendem.  Cuidado também com o estudo exagerado! É um curso de 5 a 8 anos; não dá para  manter o ritmo de estudo para vestibular por todo esse tempo. Entenda e respeite seus  limites, não exija tanto de si.  

Há certos períodos (semestres) em que poderá sentir-se à vontade para desistir de  uma disciplina em que esteja matriculade, deixando para completá-la posteriormente.  Quando você fizer isso há a possibilidade de desistir da disciplina, desmatriculando-se  oficialmente. Mas há pessoas que simplesmente deixam de cursar a disciplina,  reprovando por nota e falta e ficando com uma nota baixa em seu histórico. Cuidado  com isso, pode ser frustrante para você no futuro. Por isso, se for desistir de cursar uma  disciplina após matriculade, sempre tente pedir a desistência e para não reprovar. 

Lembre-se de que o período de graduação é muito grande, você pode mudar de  ideia a qualquer momento sobre o que pretende fazer no futuro!


\section{Pra quê estudar isso?}

O ensino médio acabou, você finalmente está livre de todas as “inutilidades”, como  química orgânica e separação silábica de verbos parnasianos, só vai ver coisas  relevantes para a profissão, e… Pimba! Uma disciplina que você acha completamente  nada a ver. Pode isso? A regra é clara, pode sim. 

Primeiro, você precisa saber que a Universidade não é um curso técnico. A ideia não é  só te dar capacitação profissional, mas sim formar pessoas melhores! Para que ume computeire precisa de contabilidade? Aparentemente para nada, mas  uma pessoa precisa ter uma noção disso, sobretudo de exatas. Outra coisa: o que  exatamente é “relevante para a sua profissão”? A computação é uma área muito vasta, e  a graduação é muito generalista para te dar base para escolher. Nosso conhecimento  pode ser aplicado em quase qualquer outra área do conhecimento com alguma dose  sucesso. Se você continuar na universidade, na pós-graduação, você só terá matérias da  sua área, já que você já sabe o suficiente pra dizer que área é essa. Mas ainda falta  muito chão até lá… Respira! 

Para quem é da Engenharia, para conseguir o CREA existem algumas matérias  obrigatórias, como Resistência dos Materiais. As coordenadorias da engenharia têm  lutado para diminuir créditos obrigatórios e aumentar eletivos – mas há matérias em  que es professores dificilmente concordariam em alterar.  

Tanto para a ciência quanto para a engenharia, o curso não é para formar simples  programadores. Vocês serão mais que isso: serão cientistas, engenheires; e isso envolve  ver coisas além de computação.

\section{Bem-vinde à academia!}

Academia é o nome que se dá à comunidade internacional de pesquisadores e  estudantes de ensino superior, atuando em todas as áreas do conhecimento. Geralmente  centrada ao redor de universidades, porém outras organizações públicas e privadas e  também empresas fazem parte. A Academia é dividida principalmente entre os pilares de  Pesquisa, Ensino e Extensão, sendo os dois primeiros mais acessíveis para nós da  graduação. A Unicamp está muito bem situada no cenário acadêmico, sendo a  responsável por parte considerável da produção científica brasileira.

\subsection{Iniciação Científica - IC}

A iniciação científica é um programa para alunes de graduação (você, no caso) ter  uma experiência acadêmica mais séria, sentir um pouco como é o clima de pesquisa.  Interessou? O que fazer? Geralmente o que se faz é conversar com professores da área  com a qual você se identifica mais (criptografia, computação teórica, processamento de  imagens, inteligência artificial etc.) e ver se estão desenvolvendo algum projeto  interessante – ou você pode propor alguma ideia sua mesmo. Depois, você começa a  estudar para redigir um projeto e o encaminha para alguma instituição de fomento à  pesquisa (CNPq ou FAPESP), pedindo uma bolsa de iniciação científica.

Além disso, você pode pegar a disciplina MC040 e posteriormente MC041, cada uma  com 12 créditos. São 24 créditos praticamente “de bandeja” para ajudá-le. Note bem as  aspas: iniciação científica geralmente consome muito tempo de estudo e dedicação, não  vá pensando que é moleza, não.

Na FEEC, você pode conseguir as matérias de iniciação (EE015 e EE016). Lá, a  iniciação científica também substitui o estágio, mas não tem equivalência com a do IC. Note que sua iniciação científica não precisa estar vinculada à computação. Como  falamos ainda neste capítulo, a Unicamp permite que você faça disciplinas de qualquer  instituto com créditos eletivos, isso pode te estimular a fazer matérias de áreas que  gosta fora da computação e nada te impede de fazer alguma iniciação científica nisso,  aproveite!

\subsection{Monitoria}

Além da iniciação científica, monitoria é uma forma muito comum de participar da  Academia, mais focada no pilar de Ensino. Não vai demorar muito para você descobrir a  importância da monitoria para as matérias. Logo no primeiro semestre, boa parte das  disciplinas possuem pessoas responsáveis por esse trabalho, que são dividas em dois  tipos: PED (Programa de Estágio Docente) e PAD (Programa de Apoio Didático).

PED é alguém da pós-graduação que é responsável por auxiliar docentes na disciplina - normalmente, criando materiais, exercícios, laboratórios e ajudando nas correções.  Já PAD é alguém da graduação e geralmente só é responsável por ajudar nos labo relatórios e tirar dúvidas, não deve incluir correção de trabalhos.

Por enquanto, o mais importante pra quem quer participar é o PAD. Todo mundo po de se inscrever, a forma e momento das inscrições varia de acordo com o instituto, mas  costuma ser no final do semestre. No IC e FEEC, recebemos um e-mail da Secretaria de  Graduação com um formulário para preencher.

O programa de PAD costuma incluir uma bolsa que fica em torno de R\$570,00. Os  pré-requisitos pra conseguir são (1) ter cursado a disciplina ou alguma equivalente, (2)  não ter reprovado na disciplina e (3) ter disponibilidade para participar das atividades,  que para o nosso curso costuma ser estar livre durante os laboratórios, (4) ter o CR  acima da média da sua turma ou ter a maior nota na matéria dentre quem se candidatou  para a bolsa.

Algumas disciplinas são mais concorridas que outras para monitoria, como é o caso  das primeiras matérias da computação. Algo interessante é que, assim como iniciação,  monitoria também tem duas disciplinas, MC050 e MC051 – ambas de 8 créditos. Mesmo  sem bolsa, é possível exercer monitoria, conseguindo os créditos. 

Mas vá com calma, como é preciso ter cursado a disciplina e ter um CR para se  candidatar pra monitoria, só vai ser possível se candidatar no final do seu segundo  semestre, já que as inscrições acontecem no fim do semestre (geralmente próximo da  matrícula). É bom se atentar que monitoria exige um certo tempo de dedicação, em  especial porque você estará ajudando outres estudantes na disciplina. É uma grande  responsabilidade e uma ótima forma de ter contato com um dos lados da Academia que  muitas vezes é esquecido: justamente o lado que lhe trouxe para cá!

\subsection{Intercâmbio}

A Unicamp é uma das universidades brasileiras que têm maior prestígio fora do país  e a VRERI-Unicamp (Vice-Reitoria Executiva de Relações Internacionais, antigo CORI), IC  e FEEC têm vários acordos bilaterais de intercâmbio. Então, para você que quer dar um  salto em algum idioma, conhecer outras culturas e buscar expandir suas oportunidades,  comece a se preparar desde já.

Muita gente fez intercâmbio pelo Ciência Sem Fronteiras, mas ele foi encerrado para  a graduação em 2017.

A França, hoje, recebe um número razoável de estudantes de computação, devido a  acordos que a Unicamp tem com os INSA e com as Écoles Centrales – e também graças a  bolsas de estudo oferecidas pela Capes e pelo governo francês. Os intercâmbios para a  França são bem mais concorridos que o Ciência sem Fronteiras, pois são na modalidade  de duplo diploma (você se forma pela Unicamp e pela instituição francesa); há um  processo seletivo envolvendo uma entrevista, e vai durar mais tempo, dois anos ou mais,  numa instituição de prestígio como a École Polytechnique, por exemplo.

No site da VRERI existem oportunidades para ir para Estados Unidos, Japão, outros  países da América Latina, Alemanha, Espanha etc, muitos com boas com incentivos que valem a pena caso você tenha um pouco de grana para se sustentar  no início. Visite sempre o site da VRERI e participe das reuniões que ela faz, pois ficar  ligade é a chave para conseguir encontrar uma boa oportunidade. Existem outras opções  de intercâmbio, como a AIESEC, que promove um intercâmbio para estágios no exterior.  Se seu interesse é mais profissional, procure se informar.

“Mas vale a pena? Poxa, vou atrasar meu curso, me deslocar da turma, vou ficar em  um país estranho, para quê? Acho que não vale a pena.”

Vamos começar pelos motivos profissionais: ter no currículo que você fala uma língua estrangeira fluentemente devido à sua imersão no país é algo muito valorizado  pelas empresas, além do fato que o pessoal do RH vai ver que você tem capacidade de  se virar sozinhe, uma vez que não é tão óbvio sair do país e recomeçar sua vida fora.  Você não atrasará tanto seu curso, pois a Unicamp conta com um sistema de  equivalências de matérias, e, se você escolher bem, pode fazer matérias que serão  convalidadas na Unicamp. Agora, o que realmente é importante: você está na faculdade,  é um bom momento de partir pro mundo! A experiência de conhecer outras culturas,  criar laços de amizade internacionais, viajar por terras desconhecidas não tem preço!  Pense que é no seu tempo de facul que terá oportunidade de fazer uma aventura destas,  não desperdice.

Se você abriu um sorriso e pensa que está preparade para sair do país, comece a  estudar! Não que ter uma boa nota seja a única forma de conseguir uma vaga em uma  bolsa de estudos, mas com certeza é a mais fácil. Busque atirar em todas as frentes,  mantenha seu CR num bom nível, busque conhecer organismos como AIESEC e procure  grupos de trabalho (no IC existem vários) que podem levar alunos ao exterior. Boa sorte!

Se você realmente se interessou, aí vão uns links com mais informações:

VRERI: \url{www.internationaloffice.unicamp.br}

AIESEC: \url{aiesec.org.br} 

Fique atente aos e-mails que você receberá do IC e da FEEC. Alguns deles são sobre  programas de intercâmbio.


\section{Acesso a artigos e revistas científicas}

Os resultados de pesquisas científicas, no Brasil e no mundo, costumam ser  publicados por meio de periódicos e conferências, os quais normalmente podem ser acessados pela internet.

No Brasil, quase todas as instituições públicas de ensino superior, como a Unicamp,  participam de um sistema conhecido como \textbf{Portal de Periódicos da Capes (\url{periodicos.capes.gov.br})}, que garante acesso a grande parte das publicações científicas  das principais editoras do mundo sem necessidade de pagar nada a mais por isso. 

Nas áreas de engenharia e de computação, quase todas as publicações relevantes são  acessíveis através deste sistema. É importante você saber, porém, que esse tipo de  acesso só é possível a partir de endereços IP da Universidade; então, se você quiser  acessar algum artigo quando estiver em casa, o ideal é usar o sistema de acesso VPN  (Virtual Private Network) disponibilizado pela Unicamp: \url{www.ccuec.unicamp.br/ccuec/servicos/acesso_remoto_vpn}. 

Através da Comunidade Acadêmica Federada (CAFe), foi disponibilizado recentemente um método de acesso remoto aos periódicos sem necessidade de usar a VPN; confira  esse link para mais informações:
\url{www-periodicos-capes-gov-br.ezl.periodicos.capes.gov.br/index.php/acesso-cafe}.

Na Unicamp, você ainda tem acesso a diversas outras publicações e e-books que não  são cobertos pelo sistema da Capes, além de alguns periódicos impressos, que podem  ser encontrados nas bibliotecas. Caso você queira buscar algo nos acervos físico e virtual  da Universidade, acesse o site do Sistema de Bibliotecas da Unicamp: \url{www.sbu.unicamp.br}.

Para uma busca mais abrangente de artigos científicos na internet, você pode usar o  Google Acadêmico (\url{scholar.google.com}). Mas atenção! Você pode encontrar artigos que  não são cobertos pelo Portal da Capes nem pela Unicamp e que exigem pagamento.

Além do Portal de Periódicos, existe também um novo modelo de publicações  científicas de acesso gratuito, chamado \textbf{\textit{open access}}. Esse modelo tem origem muito  próxima do movimento pelo software livre. Publicações feitas nesse sistema são  acessíveis a qualquer momento, de qualquer IP e sem qualquer custo. Alguns exemplos  de grandes repositórios e editoras \textit{open access} são:

\textbf{SciELO: \url{www.scielo.org}}

\textbf{PLOS: \url{www.plos.org}}

\textbf{arXiv: \url{www.arxiv.org}}

\textbf{PMC: \url{ncbi.nlm.nih.gov/pmc}}

A rede \textbf{SciELO} é onde a maior parte dos artigos em português é publicada. O acervo  \textbf{PMC} é de publicações da área biomédica.

Existe uma grande circulação de pirataria de artigos científicos na internet. Apesar  dessa prática trazer muitas coisas boas, como:

\begin{itemize}
    \item Possibilitar muitas pesquisas que não aconteceriam se pesquisadores tivessem que  pagar valores absurdos pelo acesso a conhecimentos que deveriam ser públicos;
    \item Ajudar alunes a estudarem;
    \item Democratizar o acesso ao conhecimento.
\end{itemize}

Ela também traz uma série de questões ruins, como por exemplo:

\begin{itemize}
    \item É uma grande vergonha, por isso não somos a favor de que vocês fiquem sabendo de  sites como \url{sci-hub.hkvisa.net}, \url{libgen.is} e \url{shadowlibraries.github.io}.
\end{itemize}

\section{Licenças de Software Gratuitas}

Estar na Unicamp nos dá acesso a várias licenças gratuitas de software proprietário  especiais para estudantes. 

O \textbf{LMS (Laboratório Microsoft)} ligado ao Instituto de Computação oferece download  gratuito de uma grande variedade de software da Microsoft, como o sistema Windows e  o Visual Studio. Para mais informações sobre como se cadastrar e baixar, acesse com seu  login do IC: \url{www.lms.ic.unicamp.br}. 

A \textbf{CTIC (Coordenadoria de Tecnologia da Informação e Comunicação)} obtém licenças  de muitos pacotes de software e disponibiliza para a comunidade acadêmica, inclusive  alunes. As aplicações de computação científica \textit{Wolfram Mathematica} e \textit{MATLAB} estão  disponíveis gratuitamente. Para mais informações, acesse: \url{www.ctic.unicamp.br/softwares}. 

A multinacional \textbf{Autodesk} oferece licenças de estudante gratuitas para muitos de  seus programas, incluindo o \textit{AutoCAD} e a aplicação de modelagem 3D \textit{Maya}. Para baixar,  cadastre-se usando seu e-mail da DAC, do IC ou da FEEC em \url{students.autodesk.com}. 

A \textbf{JetBrains}, empresa que produz IDEs muito boas, aprova em menos de 5 minutos os  pedidos para conta estudantil que possuem e-mail da Unicamp, o que permite o uso de  programas pagos de graça. Cadastre o e-mail do IC em \url{www.jetbrains.com/student}.

\subsection{Conta educacional no GitHub}

Além de ter acesso a alguns softwares, estudar na Unicamp também nos dá acesso a contas especiais em alguns sites, como o Github. Não se assuste com esse nome estranho, em breve você descobrirá a magia que é um sistema de controle de versão e que o git é um dos melhores deles. Além da ferramenta git por si só ser muito útil, existe o \textit{GitHub}, que disponibiliza repositórios remotos, além de uma boa forma de contribuir para projetos existentes.

Um pequeno problema do \textit{GitHub} é que em sua conta gratuita, só podemos criar  repositórios públicos, o que não combina muito com trabalhos de faculdade. O \textit{GitHub}  disponibiliza, no entanto, um \textit{Student Developer Pack}: dentre os diversos benefícios,  temos infinitos repositórios privados, descontos em serviços como a AWS, acesso às  ferramentas da \textit{Unreal Engine}, etc.

Para conseguir isso, basta utilizar seu e-mail acadêmico ao se inscrever pelo link \textbf{\url{education.github.com}}.

\subsection{Google Apps for Education}

Em 2016, a Unicamp firmou um acordo com a \textbf{Google} para disponibilizar o Google  Apps for Education para estudantes, docentes e funcionáries. Logo no começo das aulas,  você vai receber uma conta que pode acessar em \textbf{\url{googleapps.unicamp.br}}; seu login será  seu RA, e sua senha, a mesma que você usa para entrar no sistema da DAC.

Nessa conta, você vai ter acesso a tudo o que uma conta normal do Google tem, mas  com armazenamento ilimitado tanto no Gmail quanto no Drive! Ou seja, podemos  armazenar tudo o que quisermos sem preocupação! Além disso, é mais um email  acadêmico no formato: \textbf{<inicial do nome><RA>@g.unicamp.br} que você pode usar a bel prazer.

Você também pode usar Drives de equipe do Google Drive da Unicamp, que é algo  como uma pasta compartilhada de espaço ilimitado (em breve será limitado...). A  duração das chamadas pelo Google Meet também não tem limite, dentre outros  benefícios.

\section{Agora eu tenho emails da Unicamp!}

\section{Engenharia de Computação:\\ Modalidades}

Aos engenheires, esta seção serve para dar uma breve explicação sobre as duas  modalidades de graduação em engenharia de computação. Mas, primeiro, “o que é uma  modalidade?” 

Uma modalidade é uma subdivisão do curso de engenharia, um enfoque específico  da sua formação. São catálogos alternativos que visam distribuir diferentes disciplinas  para a mesma formação, ou seja, possuem ênfases em áreas diferentes para se formar  em engenharia de computação. Há duas modalidades: AA (Sistemas de Computação) e  AB (Sistemas e Processos Industriais). 

Uma das características da Unicamp é dar ao aluno uma formação MUITO  generalista, então todes es alunes, independente da modalidade – ou sendo até mesmo  da ciência –, estarão capacitados a atuar em qualquer área da computação. Assim, a  escolha da modalidade é apenas um modo de dar a alunes a oportunidade de se  aprofundar em assuntos que lhe interessem. É importante ressaltar que ambas as  modalidades têm uma grande parte de disciplinas em comum, então a formação básica  é essencialmente a mesma. Perto do período de escolha, o CACo lhe dará mais detalhes no \textbf{conjunto de palestras sobre as modalidades AA e AB}, que são realizadas no segundo  semestre do ano, para sanar todas as suas dúvidas e te auxiliar na escolha. 

\subsection{Modalidade AA: Sistemas de Computação}

Também conhecida como Azóide, é a que mais se assemelha à ciência da computação por focar mais nas áreas de análise e projeto de algoritmos, matemática discreta e  arquitetura de computadores. As aulas serão ministradas majoritariamente pelo IC.

\subsection{Modalidade AB: Sistemas e Processos Industriais}

Também conhecida como Bzóide, um pouco mais próxima da engenharia elétrica por  conta da conexão com análise de sinais, sistemas autônomos e embarcados. As aulas  serão oferecidas majoritariamente na FEEC. 

“E como escolho a modalidade?” A não ser que você tenha concluído 110 créditos  dentre as matérias da modalidade desejada com antecedência, você só fará escolherá a  modalidade desejada no final do quarto semestre, com um pedido na DAC durante o  período de matrícula ou alteração de matrícula. 

Ah, como a Unicamp te dá a oportunidade de escolher livremente suas matérias, você pode mesclar matérias tanto oferecidas para Azóides quanto para Bzóides, verificando  com veteranes como foram os oferecimentos de cada lado e escolhendo o que mais lhe  agrada; mas tome cuidado com a questão da equivalência de disciplinas: às vezes,  alguma matéria Azóide não completa a matéria Bzóide e vice-versa! 


\section{Melhores banheiros}

Uma das maiores necessidades do ser humano pode ser potencializada se for realizada num banheiro decente. Portanto, é muito importante que você saiba onde ir!  Alguns dos melhores banheiros da Unicamp são:

\textbf{IC-3}: Geralmente estão limpos e utilizáveis, mas de vez em quando têm um cheiro  ruim. E em dia de chuva ficam imundos. Exceto nos finais de semana, sempre possui  papel higiênico: é uma boa pedida na hora do apuro.

\textbf{IC-3,5}: Os banheiros do térreo não são flor que se cheire, porém, devido à localização  um pouco mais distante, pode ser suficiente se você prefere um pouco mais de  privacidade. O grande segredo são os banheiros da pós-graduação, localizados no  primeiro andar. Estes são ainda menos frequentados e a qualidade é extremamente  superior aos do andar inferior. Além disso, a torneira é excelente pra escovar os dentes  depois do almoço. Nesses casos, subir algumas escadas compensa bastante.

Observação: se o bebedouro perto dos laboratórios estiver com um gosto na água, o  bebedouro do primeiro andar do IC-3,5 tem uma vazão excelente e é um ótimo  substituto. Perfeito para encher garrafinhas de água também.

\textbf{IC-2}: Quase sempre estão limpos e utilizáveis e tem um odor melhor que os do IC-3.  Só precisa tomar cuidado pois às vezes falta papel higiênico.

\textbf{FEEC}: Possui excelentes banheiros escondidos por lá, principalmente após as refor mas de 2013. Procure bem! 

\textbf{PB}: Os banheiros do segundo e do terceiro andar do Pavilhão Básico são muito  melhores que os do térreo (especialmente os do terceiro andar, por quase não serem  usados). Só tome cuidado, porque às vezes não tem papel higiênico. 

\textbf{FE}: A Faculdade de Educação tem poucos banheiros masculinos. Estão entre os  melhores da Unicamp pelo pouco uso. 

\textbf{CB}: Estes banheiros ficam escondidos próximo às escadas do CB (no térreo). Se você  tiver sorte de chegar bem após a limpeza, o banheiro estará em excelentes condições.  Porém, na maior parte do tempo, eles ficam bem sujinhos. 

\textbf{DEQ}: Departamento de Eletrônica Quântica, no IFGW. Dizem que ninguém os usa. 

\textbf{DRCC}: Departamento de Raios Cósmicos e Cronologia, no IFGW. Um dos melhores  banheiros existentes na Unicamp (se não O melhor). Dizem que ninguém os usa. 

\textbf{DFA}: Departamento de Física Aplicada, no IFGW. Os dois andares do departamento  tem banheiros bons e utilizáveis, mas algumas vezes falta papel higiênico. 

\textbf{IMECC}: Todos os três departamentos (andares) do IMECC têm banheiros bons e  utilizáveis, mas vez ou outra falta papel higiênico. Evite os banheiros do térreo,  geralmente estão bem nojentos.


\section{Eduroam: conectando no wi-fi do campus}

O eduroam é nosso wi-fi. Parte de um projeto internacional e funciona automaticamente em várias universidades, uma vez configurado. Cuidado com seus torrents!

As instruções são diferentes para cada dispositivo, então basta se conectar na rede eduroam-config e seguir o \textbf{passo a passo disponível no site do CCUEC: \url{https://www.ccuec.unicamp.br/ccuec/servicos/rede_sem_fio_eduroam}}.

\section{A atual situação da ciência no país}

Como você viu neste manual, a universidade pública tem sido atacada por todos  os lados. Na seção de Convivência, você verá os aspectos ideológicos desses ataques; na  seção Como Funciona a Universidade, verá que os ataques orçamentários tiram uma grana  que faz falta para a pesquisa e a permanência, além do projeto de entrega total e  terceirização do ensino superior através dos “Future-se” da vida. Em 2019, nossos  pesquisadores não sabiam mais se teriam bolsa - ou universidade pública - até o fim do  ano. 

Não é fácil. \textbf{A ciência e tecnologia de um povo só avançam com base em políticas  públicas sérias e investimento pesado – que retorna para o povo}. Mas esse não parece  ser o plano para nosso país, em que as obras do acelerador síncrotron quase foram  suspensas, o Museu Nacional virou cinzas e a Amazônia (ainda) é uma floresta onde  cada vez mais se pode garimpar, desmatar, minerar e matar. 

Na Unicamp, e especialmente na Computação, você talvez tenha a noção falsa de  que estamos protegidos desses ataques. Tem bolsa FAPESP, tem interna, tem de empresa,  intercâmbio… por enquanto. Nada disso é garantido. 

Não se engane, é no olho do furacão que o vento parece calmo. A universidade é  grande, mas é só uma parte de todo um ecossistema de geração de ciência que é tão  frágil quanto um decreto, lei ou fake news – e, por isso, precisa ser defendido, quase  sempre com unhas e dentes. 

Nossa única “sorte” é que, unidos e mobilizados, vencemos: as grandes marchas  pela Educação em maio de 2019 conseguiram reverter cortes e demonstrar ao governo  que não aguentaremos o desmonte calades e que temos apoio da sociedade! 


\section{Como funciona a universidade}

\subsection{Instituição}

A universidade é uma autarquia pública - o que significa que é um órgão público do  estado, mas muito mais independente de governos (que trocam a cada 4 anos). Portanto, tem sua própria autonomia para orçamento e contratações, realizar obras, etc.

O que garante essa independência é a chamada Autonomia Universitária, garantida pelaU Constituição de 1988, mas que deriva da própria origem da "universidade" como
instrumento de avanço da sociedade como um todo. O famoso tripé que coordena (ou
deve coordenar) o trabalho da Universidade é o seguinte:

\textbf{Ensino}: formar cientistas de diversas áreas do conhecimento;

\textbf{Pesquisa}: orientar o trabalho desses profissionais à procura de novo conhecimento;

\textbf{Extensão}: retorno direto da universidade à comunidade que a sustenta, com iniciativas de impacto direto para os cidadãos e toda a sociedade. Hoje, há muitos cursos pagos (e caros) distorcendo a Extensão.

\subsection{Órgãos}

\subsection{Quem escolhe e Reitore?}

A cada 3 anos, há uma votação para escolher e reitore, o que faz muitos acharem que  é uma eleição. Contudo, é apenas uma consulta à comunidade, em que o voto des  professores tem maior peso (70\%) que o voto de funcionáries e alunes. Es nomes des 3  principais candidates mais votades nesta consulta são enviados ao Governador do  Estado, que escolhe oficialmente e reitore.

\subsection{Orçamento}

Pela Constituição, Educação é um direito da sociedade, provida pelo estado, através  dos (caros) impostos. Pela lei, as universidades estaduais paulistas recebem uma fatia do  ICMS, um imposto sobre consumo pago sobretudo por trabalhadores comuns; 
Os governos têm calculado o repasse após descontar despesas do ICMS ou ignorar  pagamentos atrasados, diminuindo o bolo investido nas universidades. Um aumento  para a Unicamp, prometido com a expansão em Limeira, não foi executado; o repasse  total para as 3 estaduais poderia chegar a 9,57\% do ICMS total. 

Há batalhas judiciais e na ALESP para assegurar os 9,57\% ou ampliá-los para 11\%,  tendo em vista o grande aumento de vagas nos últimos vários anos e os gastos fixos da  universidade, que atualmente dependem de um imposto variável. A resposta sempre é  "não há dinheiro". O governo federal apresentou às suas universidades o Future-se,  programa pelo qual as universidades entregam sua administração a empresas gestoras  (OSs) e investidoras de fundos. 

O modelo é duramente criticado, tanto pelos casos de corrupção envolvendo OSs em  concessões de hospitais públicos quanto pela extinção, na prática, da Autonomia  Universitária, atacada pelo ministro da Educação. 

Na esteira de outras universidades, a Unicamp passou em votação no CONSU a  Política de Inovação Institucional, apelidado de "Inove-se", além de outros pacotes que  vão assegurar um orçamento "extra". A regulamentação do Fundo Patrimonial da  Unicamp, por exemplo, prevê que um conselho composto por investidores e membros da  administração universitária gerencie recursos privados para prédios, bolsas e  contratações para fins específicos escolhidos pelos gestores do fundo.

Na prática, isso significa redirecionar a estrutura universitária para objetivos que se  aproximam de interesses do mercado e podem se distanciar da pesquisa de interesse  público (que não necessariamente dá lucro).

O ex-governador do estado João Doria promoveu, em 2019, na ALESP, uma CPI das  Universidades para investigar "doutrinação ideológica", com deputados aliados ao  governo federal. Reitores, ex-reitores e funcionáries foram chamados para depor. O  relatório final da CPI termina recomendando a cobrança de mensalidades. Em 2020, foi  aprovada a cobrança de mensalidades na pós-graduação lato sensu após entendimento  positivo do Superior Tribunal Federal quanto à legalidade disso em universidade  pública. 

\subsection{Como desmontar a educação para gastar menos}

Por acaso o governo está te pressionando pra cortar gastos e você não consegue dizer não? Calma, a gente tem uma solução bem simples pra isso e que quase não tem  contraindicações: sucateie a educação da sua universidade! Não custa quase nada, só a  produção científica, a educação de qualidade, a saúde mental e preparo acadêmico de  estudantes, a posição de referência científica e a autonomia universitária!



\chapter{Convivendo na Unicam}
\lipsum[1-5]

\section{whatever}
\lipsum[1-5]


\chapter{outro capitulo}
\lipsum[1-5]

%\chapter{Convivendo na Unicamp\\ {\color{white}. } \hspace{5mm} {\normalsize e no mundo}}
\chapter{Convivendo na Unicamp}
\lipsum[1-5]

\setcounter{section}{17}
\section{Mais secao (numero 17)}

\lipsum[1-5]

\section{E outra secao}

\lipsum[1-5]

\section{mais outra secao}

\lipsum[1-5]

\section{Eu amo secao}
\setcounter{subsection}{3}
\subsection{Algumas formas de evitar e combater o \\ machismo no dia a dia da universidade}

\lipsum[1-5]

%\chapter{Vivendo em Barão \\ {\color{white}. } \hspace{5mm} {\normalsize Barão Geraldo ou BG}}
\chapter{Vivendo em Barão Geraldo}
\lipsum[1-5]

\chapter{Além da Graduação}
\lipsum[1-5]

\chapter{aba}
\section{aoba}
oi

terminal \faTerminal, linux \faLinux, sapo \faFrog

github \faGithub, linkedin \faLinkedin, gitlab \faGitlab, email \faEnvelope, telegram \faPaperPlane[regular] \faTelegram\ \faTelegramPlane, whats \faWhatsapp

map marker \faMapMarker*, restaurantes \faUtensils, cafe \faCoffee, ra \faIdCard[regular] \faAddressCard[regular], casa \faHome

\begin{list}{\faTerminal}{}  
\item item A bla bla bla
\item item B blablabla
\end{list}

\section{Testes de links}


\instagram{meuinstagram}

\telegram{meugrupotelegram}

\discord{meugrupodiscord}

%\pagebreak

\lipsum[1-5]

\begin{tags}
    \github{meuusergithub} \sep \gitlab{meuusergitlab}    
\end{tags}

\end{document}
